% page 222 of the facsimile (image p-221 of the scan)
% block 165.1 x 246.3 mm ; left margin 8.0 mm
\pgc{-12.700mm}{0.000mm}{
\l{\cel{4}214}
\l{}
\l{\sou{heterogenezo}. (\sou{biol}.) Produkto di ento vivoza, da enti vi-}
\l{\cel{3}voza pre-existanta, ma di qui la naturo diferas.}
\l{}
\l{\sou{hete}rokromosomo. (\sou{biol.}) Kromosomo qua diferas de la cetera}
\l{\cel{3}kromosomi normala di la spermatozoido per sua formo e sua}
\l{\cel{3}dimensioni (plugrandeso o plumikreso).}
\l{}
\l{\sou{heteromorfoso}. (\sou{biol.}) Regenero anormala di parto di la korpo.}
\l{}
\l{\sou{heterosexuala}. Di qua la sexui esas nesama.}
\l{}
\l{\sou{heterostila}. (\sou{bot.)}\cel{2}Di qua la stili esas longa ne-egale.}
\l{}
\l{\sou{heterotipa}. (\sou{teratolo}gio). Monstro duopla en qua un ek la un-}
\l{\cel{3}aji, acesora e parazita, nekomplete developita, esas en-}
\l{\cel{3}plantacita en la parto avana di la korpo di la unajo precipua.}
\l{}
\l{\sou{heuristiko}. Parto di la historio-cienco qua traktas la serchado}
\l{\cel{3}di la dokumenti.}
\l{}
\l{\sou{hexaedra}. (\sou{geom.}) Sis-edra. - DEF.}
\l{}
\l{\sou{hexagono}. (\sou{geom.}) Figuro geometriala kun sis anguli e sis}
\l{\cel{3}lateri. - DEF.}
\l{}
\l{\sou{hexagramo}. I. (\sou{filoz.})\cel{2}Singla ek la 64 divino-kombinuri di}
\l{\cel{3}la Chiniani, kombinuri ek du de la 18 trigrami dispozita}
\l{\cel{3}vento-roze, e di qui singla konsistas ye la kombinuro ek}
\l{\cel{3}tri linei. - II. Asemblajo ek sis literi.}
\l{}
\l{\sou{hexametro}. (\sou{metriko}).Sis-peda verso. - DEFIRS.}
\l{}
\l{\sou{hezitar}. (\sou{netrans., pri, pro}) Ne-decidar pri la adopto di se-}
\l{\cel{3}lekto, di rezolvo. - DEFI.}
\l{}
\l{\sou{hiacinto}. I. (\sou{bot.}) Planto ek la familio \textquotedbl{}liliacei\textquotedbl{}. L.\cel{1}\sou{hya}-}
\l{\cel{3}\sou{cinthus}. - II. (\sou{mineral.}) Lapido, varietato di topazo.- DEFIRS}
\l{}
\l{\sou{hialoplasmo.}\cel{2}(\sou{biol.}) Parto homogena, amorfa e fluida di pro-}
\l{\cel{3}toplasmo.}
\l{}
\l{\sou{hiato}. I. (\sou{gram.}) Son-uno quan produktas la renkontro di du}
\l{\cel{3}vokali, ica finanta vorto, ita komencanta la vorto sequanta.}
\l{\cel{3}- II. (\sou{metaf.}) Nekontinueso inter la ceni di teatrajo, in-}
\l{\cel{3}ter la gradi di genealogio-arboro, e c. - DEFIRS.}
\l{}
\l{\sou{hibernar}. (\sou{netrans.}) (\sou{pri ula animali}).Vivar dum vintro, par-}
\l{\cel{3}torporante. - EF.}
\l{}
\l{\sou{hibrida}. - I. (planto od animalo) qua devenas de du speci. -}
\l{\cel{3}II. (\sou{gram}.) (Vorto) qua konsistas ye elementi, di qui ica}
\l{\cel{3}apartenas a linguo, ita a linguo diferanta. - DEFIRS.}
\l{}
\l{hido.\cel{2}(kemio). Korpo simpla, gaso, di qua la kombino kun}
\l{\cel{3}oxo produktas aquo. - Simbolo kemiala : H. - DEFIRS.}
}
